\bookchapter{Capítulo 19: Diagnóstico de fallos y autopreparación}

\chapterintro{Carga al menos 30 minutos}

\booksection{: muchas veces solo es que la batería se agoto por completo. Conecta el cargador original, espera al menos 30 minutos y prueba encender. Algunos móviles necesitan cargar hasta cierto nivel antes de responder tras agotarse.}

Reinicio forzado

\booksubsection{: para iPhone con pantalla completa: pulsa volumen + una vez, volumen - una vez, luego mantén pulsado el botón lateral hasta que aparezca el logo de Apple (pueden ser 10-15 segundos). Para Android generico: mantén el botón de encendido 10+ segundos hasta vibrar o encender.}

\begin{enumerate}
  \item \textbf{Reinicio en bucle en la pantalla de inicio}: esto se llama "manzana blanca" (iPhone) o "atascado en el logo" (Android). Si el reinicio forzado no funciona, puede que necesites entrar en modo de recuperación y repararlo con el ordenador. Conecta el iPhone al ordenador, abre iTunes/Finder, entra en modo de recuperación y selecciona "Actualizar" (no "Restaurar"): reinstalará el sistema pero conservará los datos. Para Android, entra en modo Fastboot o Recovery; las combinaciones de teclas varían según la marca: busca tu modelo + "entrar modo recuperación".
  \item \textbf{Si nada de lo anterior funciona}: entonces solo queda acudir a reparación profesional. Si hiciste copia de seguridad de los datos, no te preocupes demasiado.
  \item \textbf{VII. Problemas de sonido: altavoz no suena, volumen de llamada bajo}El altavoz no emite sonido
  \item \textbf{: confirma primero que no estas en modo silencio (interruptor lateral del iPhone), modo No molestar, o con el volumen al mínimo. Suena obvio, pero esto representa el 40\% de los casos. Luego revisa si el Bluetooth esta conectado a algún auricular o altavoz.}La otra persona no te oye en la llamada
\end{enumerate}

\textbf{: revisa si el agujero del micrófono esta tapado — el orificio inferior del móvil se tapa fácilmente con polvo. Limpialo suavemente con un cepillo de cerdas blandas. No uses aguja ni mondadientes; danaras la malla anti-polvo.}Sonido ronco/distorsionado

\booksection{: si es por agua en el altavoz — NO uses secador de pelo con aire caliente! Absorbe la humedad superficial con papel deyja secar naturalmente. Se dice que Apps de drenaje acustico (reproducen sonidos graves para expulsar el agua) pueden ayudar; pruebalos.}

VIII. Busca reparación oficial y evita las trampas

\begin{enumerate}
  \item \textbf{Si probaste todo lo anterior sin éxito, toca buscar reparación profesional. Pero el mercado de reparaciiones es un laberinto; aquí van mis experiencias y consejos:}Busca autentico oficial, no "parece oficial"
  \item \textbf{: ve al sitio web de la marca (Apple.com.cn, Huawei.com.cn, etc.) y busca la lista de "proveedores autorizados". Las tiendas no autorizadas con logo de Apple en la calle pueden usar repuestos no originales; tras la reparación perdes la resistencia al agua.}Cosas obligatorias antes de enviar a reparar
  \item \textbf{: haz copia de seguridad de TODOS los datos (no cuentes con que no se pierdan durante la reparación), cierra todas las cuentas importantes (Apple ID, WeChat, Alipay), desactiva "Buscar mi móvil". Algunos puntos de reparación te pediran la contrasena de bloqueo; esto tiene riesgos de privacidad — mejor cierra las cuentas tu mismo antes de entregarlo.}Trampas al cambiar pantalla
  \item \textbf{: pantallas no oficiales se dividen en "cristal original reciclado", "cristal no original", "pantalla ensamblada" (todo no original). Los precios varian mucho, pero también la calidad y durabilidad. Si el móvil esta en garantía, sigue el camino oficial sin atajos.}Trampas al cambiar batería
  \item \textbf{: baterías no originales pueden tener capacidad falsificada y degradarse en 6 meses más que una original en 2 anos. El precio oficial de Apple para cambiar batería no es mucho mayor que el no oficial; recomiendo ir por la via oficial.}Trampas en recuperación de datos
\end{enumerate}

\booksection{: tras caer al agua o golpearse, si en reparación dicen “la recuperación de datos cuesta extra”, ten cuidado — algúnos comercios sin escrupulos exageran el dano para cobrar más. La recuperación de datos legítima cotiza primero y luego opera; si no se recupera, no cobra.}

Truco universal numero 1: reiniciar el móvil

\begin{itemize}
  \item \textbf{Cualquier problema que tengas — lentitud, cierres inesperados, sin red, sonido anormal — reinicia primero el móvil. Suena a excusa, pero por mi experiencia, al menos el 60\% de los pequenos problemas se resuelven tras reiniciar. El móvil es como un ordenador: con el tiempo acumula cosas raras en memoria; reiniciar es empezar de cero. Asi que la próxima vez que tu móvil falle, antes de buscar "que hacer", reinicia y mira.}Memoriza ya estas operaciones salvavidas
  \item \textbf{Reinicio forzado}: para iPhone con pantalla completa: pulsa volumen + volumen - mantén pulsado el botón lateral hasta que aparezca el logo. Para Android: mantén el botón de encendido 10-15 segundos hasta que vibre.
  \item \textbf{Limpiar puerto de carga}: busca una pequeña herramienta (mondadientes o aguja de expulsión de SIM) y limpia suavemente el polvo del puerto de carga.
  \item \textbf{Reconectar Wi-Fi}: ve a Ajustes de Wi-Fi, "olvida" tu red actual y vuelve a introducir la contraseña para conectar.
\end{itemize}

\booksection{Activar/desactivar modo avión}

: desliza el centro de control, activa el modo avión, cuenta mentalmente hasta 5, desactívalo y comprueba si la red se ha recuperado.

\begin{enumerate}
  \item \textbf{Haz capturas de pantalla de estos pasos o scribelos en notas. Cuando tengas problemas, no entres en panico buscando; simplemente abre tus notas.}Capítulo 20: Modo desarrollador y permisos avanzados
  \item \textbf{Este capítulo es para quienes quieren explorar las "funciones ocultas" del móvil. Modo desarrollador, permisos Root, bloqueo BL... estos conceptos suenan muy geek, pero en realidad no estan lejos del usuario común.}I. Modo desarrollador: desbloquear ajustes ocultos
  \item \textbf{El modo desarrollador es un menu oculto integrado en el sistema, con muchas funciones que los usuarios normales no necesitan pero los desarrolladores si.}Cómo activarlo:
  \item \textbf{Android}: ve a Ajustes > Sobre el móvil y pulsa "Número de versión" 7 veces seguidas hasta que aparezca el mensaje "Ya estás en modo desarrollador".
\end{enumerate}

\booksection{HarmonyOS}

\begin{itemize}
  \item \textbf{: Ajustes > Sobre el móvil > pulsa "Versión de HarmonyOS" 7 veces seguidas.}iOS
  \item \textbf{: los usuarios comunes no pueden activar el modo desarrollador completo; se necesita una cuenta de Apple Developer de pago.}Funciones habituales:
  \item \textbf{Depuracion USB}: conectar al ordenador para operaciones avanzadas como flasheo, recuperación de datos.
  \item \textbf{Ajuste de velocidad de animación}: reducir o desactivar animaciones para que el móvil se sienta más rápido.
\end{itemize}

\booksection{Mostrar toques en pantalla}

: muestra puntos de toque en la pantalla, ideal para grabar tutoríales.

\begin{itemize}
  \item \textbf{Limitar procesos en segúndo plano}: limita el numero de Apps en segúndo plano para ahorrar batería.
  \item \textbf{Ventajas de Root:}: mejora la fluidez en algúnos casos.
  \item \textbf{Notas importantes}Los ajustes del modo desarrollador pueden afectar la estabilidad del sistema; no modifiques opciones que no entiendas. Tras activar la Depuración USB, ten cuidado al conectar al ordenador: no conectes a ordenadores en los que no confíes.
  \item \textbf{II. Bloqueo BL (Bootloader): primera puerta del flasheo}El Bootloader es el programa que carga el sistema operativo al iniciar. El bloqueo BL es un candado que ponen los fabricantes por seguridad para evitar que los usuarios instalen sistemas no oficiales.
\end{itemize}

\booksection{Funciones del bloqueo BL:}

\begin{itemize}
  \item \textbf{Protección de seguridad}: evita que software malicioso flashee el sistema.
  \item \textbf{Anti-robo}: incluso si el móvil es flasheado, no se puede usar sin la cuenta y contraseña originales.
  \item \textbf{Limitación de garantía}: desbloquear el BL suele significar perder la garantía oficial.
\end{itemize}

\booksection{Políticas de bloqueo BL por marca:}

Xiaomi

\begin{itemize}
  \item \textbf{: se puede desbloquear; requiere vincular la cuenta y un período de espera (7-30 días).}Huawei
  \item \textbf{: desde 2018 dejó de proporcionar códigos de desbloqueo; básicamente imposible desbloquear oficialmente.}OnePlus
  \item \textbf{: se puede desbloquear; política relativamente flexible.}OPPO/vivo
  \item \textbf{: algúnos modelos se pueden desbloquear; requiere solicitud.}Samsung
  \item \textbf{: se puede desbloquear, pero activa el fusible Knox y pierdes Samsung Pay y otras funciones.}Apple
\end{itemize}

\begin{tcbwarn}{: los dispositivos iOS no tienen el concepto tradicional de bloqueo BL, pero tienen mecanismos de seguridad similares.}
III. Permisos Root: el “superadministrador” de Android

\end{tcbwarn}

\begin{tcbpractice}{Los permisos Root son equivalentes al Administrador de Windows; es el nivel máximo de permisos en Android. Tras obtener permisos Root, puedes modificar archivos del sistema, desinstalar Apps preinstaladas y usar Apps avanzadas que requieren Root.}
\begin{enumerate}
  \item \textbf{Ventajas de Root:}Desinstalar Apps preinstaladas que no se pueden desinstalar normalmente.
  \item \textbf{Usar Apps con funciones avanzadas (como Titanium Backup, Greenify).}Modificar la interfaz y funciones del sistema.
  \item \textbf{Instalar Recovery de terceros y ROMs personalizadas.}Riesgos de Root:
  \item \textbf{Pérdida de garantía}: la mayoría de fabricantes no garantizan móviles con Root.
  \item Riesgos de seguridad
\end{enumerate}

\end{tcbpractice}
